\section*{Frontmatter reader architecture}
\addcontentsline{toc}{section}{Frontmatter reader architecture}
\begin{quote}\small
\noindent\textbf{Purpose.} This page separates the human reading route from the LLM/audit route. It is an orientation layer only; it does not replace any definition, theorem, proof, audit ledger, or historical witness later in the document.
\end{quote}

\begin{table}[H]
\centering
\small
\begin{tabular}{>{\raggedright\arraybackslash}p{0.24\textwidth} >{\raggedright\arraybackslash}p{0.31\textwidth} >{\raggedright\arraybackslash}p{0.31\textwidth}}
\toprule
\textbf{Layer} & \textbf{Human reading route} & \textbf{LLM/audit route}\\
\midrule
Main idea & Read the Sphaera, HC, MM, and OP interfaces as the conceptual spine. & Track named definitions, route tags, residual triples, and non-certificate markers.\\
QAM/MML/MMA & Treat QAM as store, MML as commands, and MMA as execution/readout semantics. & Use the finite object bridge, instruction alphabet, and deterministic guarded transitions as machine-readable handles.\\
Tool layer & Read repeated proof endings through compact tool contracts. & Preserve first occurrences, hypotheses, guard pointers, witness pointers, and theorematic status.\\
Sphaera array & Read several Sphaera objects as a guarded mirror/stripe/parity working array. & Track component stores, cross-sphere guards, residual/parity witnesses, snapshots, and scrub outputs.\\
Pre-v4 QAM shift & Read tag 0 as a future reserved empty marker, not as a current migration. & Track the Hilbert-shift plan, conservative reindexing guard, mixed-readout packages, and optional binary masks.\\
Two-sided tensor grammar & Read it as a notation guide, not as a new tensor calculus. & Track front/readout-side indices, back/structure-side indices, and the rule that same-variance pairings need explicit comparison or selector tensors.\\
Historical witness & Use it as genealogy and reconstruction support. & Preserve it as non-destructive proof memory; do not read inherited baselines as current OP closure.\\
\bottomrule
\end{tabular}
\caption{Frontmatter separation of human readability and LLM/audit precision.}
\label{tab:frontmatter-human-llm-architecture}
\end{table}

\subsection*{Version and DOI record for v3.35.0}
\addcontentsline{toc}{subsection}{Version and DOI record for v3.35.0}
\label{subsec:v3350-version-doi-record}

\begin{center}
\small
\begin{tabular}{>{\raggedright\arraybackslash}p{0.24\textwidth} >{\raggedright\arraybackslash}p{0.24\textwidth} >{\raggedright\arraybackslash}p{0.40\textwidth}}
\toprule
\textbf{Record} & \textbf{Date / status} & \textbf{DOI / note}\\
\midrule
Concept DOI & ongoing pointer & \href{https://doi.org/10.5281/zenodo.19210728}{10.5281/zenodo.19210728}\\
Prior public version & v3.34.0, 2026-04-25 & \href{https://doi.org/10.5281/zenodo.19746821}{10.5281/zenodo.19746821}\\
Current public version & v3.35.0, 2026-04-25 & \href{https://doi.org/10.5281/zenodo.19750674}{10.5281/zenodo.19750674}\\
\bottomrule
\end{tabular}
\end{center}

\noindent\textbf{DOI rule.} The concept DOI remains the stable title-page pointer of the Companion line. The version-specific DOI is release metadata only; it is not used as a mathematical index, subscript, theorem label, or object name.

\paragraph{Current local working status.}
The active public baseline is v3.34.0. The present v3.35.0 release is a pre-v4 preparation pass. It keeps the split-bundle architecture and adds only preparatory QAM material: a short/extended abstract architecture, a Hilbert-hotel reindexing plan for later v4, guarded mixed QAM readouts as derived packages, an optional binary mask registry for readout classification, and a conservative two-sided structural tensor grammar for local readout/action/comparison notation. It does not globally migrate existing QAM tuple tags and does not globally rewrite the Companion into the new notation.

\paragraph{No-overclaim guard.}
All frontmatter descriptions in this page are routing aids. They do not assert OP closure, residual vanishing, technical RAID reconstruction, or Turing completeness. Any later upgrade to those claims must be discharged by a separate explicit theorem.


\paragraph{r05 public split-bundle release note.}
The r05 layer is a packaging and readability finalization pass. It keeps the active core as the governing current proof line and keeps the historical witness in a separate PDF of the same bundle. The separation is archival and typographical: it does not delete historical material, does not promote older historical statements to current theorematic status, and does not add a closure certificate.

\clearpage
\tableofcontents
\clearpage

